\section{Conclusion}
\label{sec:conclusion}

Should a language model be asked what to order? The question is badly posed,
because it silently bundles three separable decisions, and separating them is what
makes any resulting gain attributable. \collie{} consults the model rarely,
restricts it to a typed claim about what changed, compiles that claim into one
point of a registered grid over a controller the operations research literature
already trusts, and lets the claim touch an order only after an anytime-valid
\eprocess{} started strictly after the claim was frozen has paid for it. The
model says what changed, once. Statistics decides whether to trust it. Operations
research decides what to do.

Two things we built are useful whether or not the method wins. One is an audited
testbed whose baseline reproduces every published order sequence bit for bit,
which is what let us establish that the published baseline score belongs to one
of the benchmark's two simulators and is unreachable from the other, so any future
comparison here has to name its harness. The other is a ledger that keeps
physical requests separate from charged copies, which makes a shared proposal
call across three activation policies a controlled ablation rather than a cost
accident. \pending{final result paragraph: which contrasts survived Holm
correction, what the frontier showed, and which kill triggers fired}
